\documentclass[12pt, letterpaper]{article}
\usepackage[utf8]{inputenc}
\usepackage[spanish]{babel}
\usepackage{graphicx}
\usepackage{amsmath}
\usepackage{amsfonts}
\usepackage{amssymb}
\usepackage{geometry}
\usepackage{setspace}
\usepackage{tocloft}
\usepackage{parskip}
\usepackage{hyperref}
\usepackage{fancyhdr}
\usepackage{array}
\usepackage{listings}
\usepackage{xcolor}
\usepackage{float}
\usepackage[style=apa, backend=biber]{biblatex}
\addbibresource{references.bib}

\geometry{
    letterpaper,
    total={170mm,257mm},
    left=30mm,
    right=30mm,
    top=25mm,
    bottom=25mm,
}
\renewcommand{\cftsecleader}{\cftdotfill{\cftdotsep}}
\setlength{\cftbeforesecskip}{1em}

%=========== ESTILO PARA EL CÓDIGO ===========
\definecolor{codegreen}{rgb}{0,0.6,0}
\definecolor{codegray}{rgb}{0.5,0.5,0.5}
\definecolor{codepurple}{rgb}{0.58,0,0.82}
\definecolor{backcolour}{rgb}{0.95,0.95,0.92}

\lstdefinestyle{sqlstyle}{
    backgroundcolor=\color{backcolour},
    commentstyle=\color{codegreen},
    keywordstyle=\color{magenta},
    numberstyle=\tiny\color{codegray},
    stringstyle=\color{codepurple},
    basicstyle=\footnotesize\ttfamily,
    breakatwhitespace=false,
    breaklines=true,
    captionpos=b,
    keepspaces=true,
    numbers=left,
    numbersep=5pt,
    showspaces=false,
    showstringspaces=false,
    showtabs=false,
    tabsize=2,
    language=SQL
}

\lstdefinestyle{dbmlstyle}{
    backgroundcolor=\color{backcolour},
    commentstyle=\color{codegreen},
    keywordstyle=\color{blue},
    stringstyle=\color{codepurple},
    basicstyle=\footnotesize\ttfamily,
    breaklines=true,
    captionpos=b,
    keepspaces=true,
    numbers=left,
    numbersep=5pt,
    showstringspaces=false,
    morekeywords={Table, integer, primary, key, increment, varchar, unique, not, null, ref, numeric}
}

\pagestyle{fancy}
\fancyhf{}
\renewcommand{\headrulewidth}{0pt}
\cfoot{\thepage}

% ═══════════════════════════════════════════════════════════════════════════════
\begin{document}

% ── PORTADA 1 ─────────────────────────────────────────────────────────────────
\begin{titlepage}
    \centering
    \vspace*{0.5cm}
    {\large \bfseries Software de implementación de funcionalidades y aplicaciones de un modelo
    matemático para el análisis y optimización de mezclas de concreto estructurales\par}
    \vspace{1.5cm}
    % Coloca el logo en docs/images/logo_ucc.png
    % \includegraphics[width=0.35\textwidth]{images/logo_ucc.png}\par
    \vspace{1.5cm}
    {\large Auxiliar de investigación\par}
    \vspace{1.5cm}
    {\large \bfseries Juan Diego Pedroza Sierra\par}
    \vspace{3cm}
    {\normalsize Universidad Cooperativa de Colombia\par}
    {\normalsize Facultad de Ingeniería\par}
    {\normalsize Pregrado de Ingeniería de Sistemas\par}
    {\normalsize Villavicencio, Colombia\par}
    {\normalsize 2026\par}
\end{titlepage}

% ── PORTADA 2 — ASESOR ────────────────────────────────────────────────────────
\newpage
\thispagestyle{empty}
\begin{center}
    \vspace*{0.5cm}
    {\large \bfseries Software de implementación de funcionalidades y aplicaciones de un modelo
    matemático para el análisis y optimización de mezclas de concreto estructurales\par}
    \vspace{2cm}
    {\large \bfseries Juan Diego Pedroza Sierra\par}
    \vspace{2cm}
    {\normalsize Asesor(es)\par}
    {\normalsize Carlos Ignacio Torres Londoño\par}
    {\normalsize Gyovany\par}
    \vspace{2cm}
    {\normalsize Trabajo de grado presentado para optar al título de Ingeniero de Sistemas\par}
    \vspace{3cm}
    {\normalsize Universidad Cooperativa de Colombia\par}
    {\normalsize Facultad de Ingeniería\par}
    {\normalsize Pregrado de Ingeniería de Sistemas\par}
    {\normalsize Villavicencio, Colombia\par}
    {\normalsize 2026\par}
\end{center}

% ── AUTORIDADES ACADÉMICAS ────────────────────────────────────────────────────
\newpage
\thispagestyle{empty}
\begin{center}
    \vspace*{2cm}
    {\large \bfseries Autoridades académicas Universidad Cooperativa de Colombia\par}
    \vspace{2cm}
\end{center}
\begin{center}
    \large
    Dra. Maritza Rondón Rangel\\
    \textit{Rectora Nacional}\\[1cm]
    Dr. César Augusto Pérez Londoño\\
    \textit{Director de la Sede}\\[1cm]
    Dr. Fabián Ignacio Ojeda\\
    \textit{Subdirector Académico de la Sede}\\[1cm]
    Dra. Adriana Carolina Urrea Ortiz\\
    \textit{Subdirectora de Desarrollo Institucional y Financiero}\\[1cm]
    Ing. María Lucrecia Ramírez Suárez\\
    \textit{Decana (E) programa de Ingenierías}\\[1cm]
    Ing. Pedro Alexander Gutiérrez Aguilera\\
    \textit{Coordinador de Investigación de la Facultad de Ingenierías}\\[1cm]
\end{center}

% ── ADVERTENCIA ───────────────────────────────────────────────────────────────
\newpage
\thispagestyle{empty}
\begin{center}
    \vspace*{3cm}
    {\large \bfseries Advertencia\par}
    \vspace{2cm}
\end{center}
\begin{flushright}
    \large
    La Universidad Cooperativa de Colombia, sede Villavicencio, no se hace\\
    responsable por los conceptos emitidos por los autores.
\end{flushright}

% ── NOTA DE ACEPTACIÓN ────────────────────────────────────────────────────────
\newpage
\thispagestyle{empty}
\begin{center}
    \vspace*{3cm}
    {\large \bfseries Nota de aceptación\par}
    \vspace{2cm}
    \rule{10cm}{0.4pt}\\[0.3cm]
    \rule{10cm}{0.4pt}\\[0.3cm]
    \rule{10cm}{0.4pt}\\[0.3cm]
    \rule{10cm}{0.4pt}\\[0.3cm]
    \rule{10cm}{0.4pt}\\[0.3cm]
    \rule{10cm}{0.4pt}\\[0.3cm]
    \rule{10cm}{0.4pt}\\[0.3cm]
    \rule{10cm}{0.4pt}\\[0.3cm]
    \rule{10cm}{0.4pt}\\[0.3cm]
    \rule{10cm}{0.4pt}\\[1.5cm]
    \rule{8cm}{0.4pt}\\
    Firma del presidente del jurado\\[1.2cm]
    \rule{8cm}{0.4pt}\\
    Firma del jurado\\[1.2cm]
    \rule{8cm}{0.4pt}\\
    Firma del jurado
\end{center}

% ── DEDICATORIA ───────────────────────────────────────────────────────────────
\newpage
\thispagestyle{empty}
\begin{flushright}
    \vspace*{3cm}
    {\large \bfseries Dedicatoria\par}
    \vspace{2cm}
    \textit{Esta página es opcional.\\
    Mencione, en pocas líneas, las personas a quienes dedica su trabajo.\\
    Ubique el texto alineado a la derecha sin justificar.}
\end{flushright}

% ── AGRADECIMIENTOS ───────────────────────────────────────────────────────────
\newpage
\thispagestyle{empty}
\begin{center}
    \vspace*{3cm}
    {\large \bfseries Agradecimientos\par}
    \vspace{2cm}
    \textit{Esta página es opcional. Mencione aquellas personas, instituciones u organizaciones\\
    que hicieron posible el desarrollo del trabajo de grado.}
\end{center}

% ── RESUMEN ───────────────────────────────────────────────────────────────────
\newpage
\pagenumbering{roman}
\setcounter{page}{1}

\section*{Resumen}
\addcontentsline{toc}{section}{Resumen}

En el departamento del Meta (Colombia), las condiciones ambientales aceleran la degradación del concreto y elevan los costos de mantenimiento. Ante la ausencia de herramientas locales para apoyar la dosificación y selección de aditivos, este proyecto desarrolla DIAPCE: una aplicación multiplataforma en Flutter dirigida exclusivamente a los estudiantes del semillero de investigación y a los participantes en prácticas de laboratorio de materiales de la Universidad Cooperativa de Colombia (UCC), sede Meta. DIAPCE no es una herramienta de uso industrial ni general; su propósito es formativo e institucional. A partir de un conjunto de datos regional y entradas como resistencia objetivo, temperatura, humedad y condiciones de exposición, el sistema sugiere una dosificación base y clasifica el uso de aditivos (P0: sin aditivo; P1--P3: impermeabilizante al 2, 3~y~4\,\%; PP1--PP3: plastificante al 0{,}2, 0{,}4 y 0{,}6\,\%), con proyecciones a 7, 14 y 28 días, en concordancia con ACI\,211.1 y normas relacionadas. La solución funciona sin conexión mediante SQLite y prioriza la trazabilidad de los ensayos y ejercicios académicos (registro de insumos, resultados y versión del modelo), facilitando prácticas reproducibles y comparables en el aula y el laboratorio. DIAPCE busca fortalecer competencias en diseño de mezcla, análisis de resultados y cultura de registro en los estudiantes de la UCC, y potencialmente contribuir a un uso más eficiente de materiales en el contexto de sus prácticas formativas.

\vspace{0.5em}
\noindent\textbf{Palabras clave:} concreto estructural, dosificación, aditivos (P0--PP3: impermeabilizante 0--4\,\% y plastificante 0{,}2--0{,}6\,\%), ACI\,211.1, Flutter, SQLite, Meta (Colombia), 7/14/28 días, laboratorios, semilleros de investigación, docencia.

% ── ABSTRACT ──────────────────────────────────────────────────────────────────
\newpage
\section*{Abstract}
\addcontentsline{toc}{section}{Abstract}

In the Meta department (Colombia), environmental conditions accelerate concrete degradation and increase maintenance costs. Given the lack of local tools to support mix design and admixture selection, this project develops DIAPCE: a cross-platform Flutter application aimed exclusively at students in the research seedbed and participants in materials laboratory practices at the Universidad Cooperativa de Colombia (UCC), Meta campus. DIAPCE is not an industrial or general-purpose tool; its purpose is purely educational and institutional. Based on a regional dataset and inputs such as target strength, temperature, humidity, and exposure conditions, the system suggests a base mix design and classifies the use of admixtures (P0: no admixture; P1-P3: waterproofing at 2, 3, and 4\,\%; PP1-PP3: plasticizer at 0.2, 0.4, and 0.6\,\%), with projections at 7, 14, and 28 days, in accordance with ACI\,211.1 and related standards. The solution works offline using SQLite and prioritizes the traceability of academic tests and exercises (recording of inputs, results, and model version), facilitating reproducible and comparable practices in the classroom and laboratory. DIAPCE seeks to strengthen skills in mix design, results analysis, and data-recording culture among UCC students, potentially contributing to a more efficient use of materials within the context of their educational practices.

\vspace{0.5em}
\noindent\textbf{Keywords:} structural concrete, mix design, admixtures (P0-PP3: waterproofing 0-4\,\% and plasticizer 0.2-0.6\,\%), ACI\,211.1, Flutter, SQLite, Meta (Colombia), 7/14/28 days, laboratories, research seedbeds, teaching.

% ── TABLAS DE CONTENIDO ───────────────────────────────────────────────────────
\newpage
\tableofcontents
\newpage
\listoftables
\addcontentsline{toc}{section}{Lista de tablas}
\newpage
\listoffigures
\addcontentsline{toc}{section}{Lista de figuras}

% ══════════════════════════════════════════════════════════════════════════════
%  CUERPO DEL DOCUMENTO
% ══════════════════════════════════════════════════════════════════════════════
\newpage
\pagenumbering{arabic}
\setcounter{page}{1}

% ══════════════════════════════════════════════════════════════════════════════
%  CAPÍTULO 1 — DEFINICIÓN DEL PROYECTO
% ══════════════════════════════════════════════════════════════════════════════
\section{Definición del Proyecto}

% ─────────────────────────────────────────────────────────────────────────────
\subsection{Descripción del problema}

En el departamento del Meta (Colombia), las condiciones ambientales agresivas ---altas temperaturas, humedad elevada y particularidades de los suelos--- inciden de manera directa en el desempeño y la durabilidad del concreto empleado en obras civiles. A pesar de ello, existe una carencia de herramientas digitales locales que apoyen, de forma didáctica y con trazabilidad, la dosificación y la selección de aditivos según el contexto regional. Esta brecha se hace evidente en los laboratorios de materiales y en los semilleros de investigación de la Universidad Cooperativa de Colombia (UCC), donde la replicabilidad de los ensayos y la sistematización de resultados son esenciales para la formación y la transferencia de conocimiento.

La ausencia de software especializado para apoyar el diseño y la verificación de mezclas adaptadas a condiciones locales puede derivar en decisiones inconsistentes, menor durabilidad y mayores costos de mantenimiento. De acuerdo con lineamientos nacionales, es prioritario impulsar la innovación tecnológica en el sector construcción para elevar la calidad y sostenibilidad de la infraestructura (Departamento Nacional de Planeación, 2018). A su vez, gremios y estudios académicos resaltan la necesidad de incorporar herramientas digitales que integren particularidades ambientales y estándares técnicos en el diseño de mezclas para optimizar recursos y asegurar la vida útil de las estructuras (CCI, 2020; González et al., 2019).

En este contexto, se identifica la necesidad de una solución \textbf{dirigida a los estudiantes del semillero de investigación y a quienes participan en las prácticas de laboratorio de materiales de la UCC}, que les permita: (i) registrar de manera estructurada las condiciones de diseño, (ii) asistir la dosificación base y la selección de aditivos con proyecciones de resistencia a distintas edades, y (iii) garantizar la trazabilidad de los ensayos y ejercicios académicos. Es importante precisar que \textbf{DIAPCE no está concebida como una herramienta de uso general ni para profesionales en ejercicio fuera del entorno universitario}; su alcance es intencionalmente formativo e institucional. DIAPCE responde a esta necesidad mediante una aplicación multiplataforma desarrollada en Flutter que opera sin conexión (SQLite), orientada exclusivamente a estudiantes del semillero y los laboratorios de la UCC para apoyar decisiones basadas en datos, estandarizar prácticas académicas y facilitar la comparación reproducible de resultados. La herramienta no sustituye el criterio del docente ni el marco normativo, sino que lo complementa con evidencia sistematizada y accesible en el contexto regional del Meta.

\subsubsection*{Pregunta de investigación}
¿Cómo por medio de la implementación de un software se pueden realizar cálculos de optimización de mezclas cementantes a partir de un modelo matemático que use propiedades físicas, mecánicas y térmicas de concreto estructural?

% ─────────────────────────────────────────────────────────────────────────────
\subsection{Objetivos}

\subsubsection{Objetivo general}
Desarrollar un software que implemente un modelo matemático para el ajuste de proporciones de materiales cementantes en mezclas de concreto estructural, orientado a los estudiantes del semillero de investigación y de los laboratorios de materiales de la UCC, sede Meta.

\subsubsection{Objetivos específicos}
\begin{enumerate}
    \item Implementar una base de datos que permita el almacenamiento de características físicas, mecánicas y/o ambientales de mezclas de concreto estructural.
    \item Desarrollar un software que implemente un modelo matemático para las mezclas de concreto estructural con una arquitectura multiplataforma y almacenamiento local.
    \item Realizar pruebas y análisis de desempeño del programa referente a criterios de calidad.
\end{enumerate}

% ─────────────────────────────────────────────────────────────────────────────
\subsection{Justificación}

\subsubsection*{Por qué surge esta propuesta}

La necesidad de DIAPCE no nace de una carencia tecnológica abstracta, sino de una realidad concreta y repetida en el laboratorio de materiales y en el semillero de investigación de la UCC, sede Meta: cada cohorte de estudiantes parte prácticamente desde cero. Los registros de ensayos anteriores no se encuentran en un formato consultable, las condiciones de cada experimento se documentan de forma dispersa y los resultados acumulados por el grupo no se traducen en un saber colectivo que oriente las decisiones de dosificación siguientes. Cuando el conocimiento generado en prácticas sucesivas no se sistematiza, la curva de aprendizaje se reinicia con cada nueva generación de estudiantes, y la ventaja de pertenecer a un semillero activo se diluye.

A esto se suma el contexto ambiental del Meta: las condiciones de temperatura y humedad del territorio no son las que suponen las tablas genéricas de dosificación derivadas de normas internacionales. Diseñar mezclas de concreto estructural en esta región exige contrastar esas referencias con evidencia experimental propia, producida y validada localmente. Sin una herramienta que centralice y haga consultable ese acervo, cada decisión de dosificación se apoya en criterios individuales difíciles de comparar y reproducir.

\subsubsection*{Qué busca el proyecto}

DIAPCE busca convertir el conocimiento experimental acumulado por el semillero en un recurso colectivo, trazable y reproducible. No se trata de reemplazar el juicio del docente ni de automatizar el diseño de mezclas; se trata de ofrecer a los estudiantes una plataforma donde los datos de sus propios ensayos ---resistencias medidas, condiciones ambientales registradas, aditivos empleados--- se conviertan en la base real de las recomendaciones de dosificación.

En términos concretos, el proyecto busca: (i) centralizar los registros de diseño y los resultados de laboratorio en un repositorio estructurado y persistente; (ii) proporcionar asistencia en la selección de proporciones y aditivos a partir de los datos experimentales del semillero, con proyecciones de resistencia a distintas edades; y (iii) establecer una práctica de registro sistemático que permita comparar resultados entre sesiones, cohortes y condiciones ambientales distintas. Todo ello dentro de un entorno de uso exclusivamente formativo, sin pretensión de reemplazar herramientas de uso profesional o industrial.

\subsubsection*{Beneficios esperados}

Para el estudiante, DIAPCE representa un punto de contacto directo entre la teoría del diseño de mezclas y la evidencia que él mismo produce en el laboratorio. La herramienta hace visible el impacto de las variables ambientales y de la selección de aditivos sobre la resistencia del concreto, reforzando competencias que van más allá de la memorización de fórmulas: análisis crítico de resultados, formulación y contraste de hipótesis, y cultura de registro.

Para el semillero como grupo, el beneficio principal es la trazabilidad: la posibilidad de recuperar, comparar y discutir resultados de ensayos anteriores como si formaran parte de una misma conversación científica continua. Esto fortalece la identidad investigativa del grupo y eleva la calidad de las discusiones académicas entre cohortes.

Para los docentes, DIAPCE ofrece un instrumento de evaluación objetiva: los registros quedan almacenados con sus parámetros de entrada, lo que permite contrastar el proceso de decisión de cada estudiante con los resultados obtenidos y con el comportamiento del conjunto de datos regional.

\subsubsection*{Componente comunitario y aporte a la comunidad académica}

El valor de DIAPCE trasciende el uso individual. Al operar sobre un repositorio compartido de datos experimentales generados por los propios estudiantes de la UCC en el contexto del Meta, la herramienta construye ---práctica a práctica--- un patrimonio de conocimiento regional. Cada ensayo registrado enriquece la base desde la que se producen las recomendaciones futuras, lo que significa que el esfuerzo de investigación de una cohorte beneficia directamente a las siguientes.

Este modelo de acumulación progresiva del conocimiento es especialmente relevante en una región donde las condiciones ambientales particulares hacen insuficientes las referencias externas y donde la producción de datos locales tiene un valor formativo e institucional difícil de sustituir. DIAPCE propone que el semillero de investigación de la UCC deje de ser un espacio donde el conocimiento se genera y se pierde por ciclos, y se convierta en un espacio donde ese conocimiento se acumula, se discute y se refina de manera continua.

% ─────────────────────────────────────────────────────────────────────────────
\subsection{Metodología}

Como metodología de trabajo se optó por usar \textbf{FDD (Feature-Driven Development)} debido a su enfoque ágil y centrado en el cliente. FDD facilita la división del trabajo en pequeñas funciones, permitiendo una detección temprana de errores y una rápida adaptación a los cambios (Lucidchart, 2024). El tipo de investigación es mixto, con un enfoque en el desarrollo de software.

\begin{figure}[H]
    \centering
    \includegraphics[width=1\textwidth]{images/metodolgia.jpg}
    \caption{Diagrama metodológico del proyecto con la estructura de la documentación y desarrollo.}
    \label{fig:metodologia}
\end{figure}

Las fases metodológicas adoptadas son:
\begin{enumerate}
    \item \textbf{Planteamiento:} Identificación del problema, definición de objetivos y revisión de antecedentes.
    \item \textbf{Diseño:} Modelado de la base de datos, arquitectura de la aplicación y especificación de requisitos funcionales.
    \item \textbf{Desarrollo:} Implementación en Flutter/Dart con almacenamiento local SQLite e integración del conjunto de datos experimental.
    \item \textbf{Pruebas:} Validación funcional de la interfaz, pruebas de calidad del modelo de predicción y revisión de criterios normativos.
\end{enumerate}

Las \textbf{variables de investigación} consideradas son:
\begin{itemize}
    \item \textbf{Variable independiente:} El modelo matemático y los parámetros de entrada (temperatura, humedad, relación a/c, aditivo).
    \item \textbf{Variable dependiente:} Los resultados generados por el software (proyecciones de resistencia a 7/14/28 días) y las métricas de calidad (MAE, RMSE, MAPE, F1).
    \item \textbf{Variables intervinientes:} Errores en los datos de entrada, desbalance de clases en el dataset.
    \item \textbf{Variables controladas:} Formatos de entrada de datos, entorno de pruebas y criterios de desempeño definidos.
\end{itemize}

% ══════════════════════════════════════════════════════════════════════════════
%  CAPÍTULO 2 — MARCO TEÓRICO
% ══════════════════════════════════════════════════════════════════════════════
\newpage
\section{Marco Teórico}

Este marco teórico presenta los fundamentos conceptuales y normativos que sustentan el desarrollo de DIAPCE. Se organiza en siete ejes: el diseño de mezclas de concreto y la ley de Abrams como base empírica, la influencia de las condiciones ambientales sobre el desempeño del concreto, los aditivos químicos y su clasificación, el desarrollo de resistencia a compresión a lo largo del tiempo, el modelado empírico adoptado por el sistema, los fundamentos tecnológicos de la solución y, finalmente, los trabajos relacionados que contextualizan la propuesta.

% ─────────────────────────────────────────────────────────────────────────────
\subsection{Diseño de mezclas de concreto estructural — Ley de Abrams y método ACI}

El diseño de mezclas de concreto es el proceso mediante el cual se determinan las proporciones de los materiales constituyentes ---cemento, agua, agregados finos y gruesos, y eventualmente aditivos--- con el fin de obtener un material que cumpla con los requisitos de resistencia, trabajabilidad y durabilidad especificados (ACI 211.1, 2016). El fundamento histórico de este proceso es la \textbf{ley de Abrams}, que establece una relación inversa entre la resistencia a compresión $f'_c$ y la relación agua-cemento (a/c). En forma log-lineal puede expresarse como:
\[
f'_c = K \cdot \left(\frac{w}{c}\right)^{-n}
\quad\text{o equivalentemente}\quad
\log f'_c = \log K - n \cdot \log\!\left(\frac{w}{c}\right),
\]
donde $K$ y $n$ son parámetros dependientes del tipo de cemento, los materiales y las condiciones de curado (Jiménez y Alvarado, 2020). A menor relación a/c, menor es la porosidad de la pasta cementante y mayor la resistencia resultante; sin embargo, la trabajabilidad de la mezcla disminuye, lo que puede comprometer la colocación y el acabado.

El método más difundido para traducir esta relación en proporciones prácticas es el del \textit{American Concrete Institute}, \textbf{ACI 211.1} (ACI 211.1, 2016), que define un procedimiento sistemático a partir de la resistencia de diseño ($f'_c$), el tamaño máximo del agregado grueso, el asentamiento requerido y la relación a/c. En Colombia, el diseño de mezclas estructurales también se rige por la \textbf{NSR-10}, cuyo título C establece los requisitos mínimos de resistencia y durabilidad. La verificación experimental se realiza conforme a la \textbf{NTC 673} y \textbf{ASTM C39} (ASTM C39, 2021), mediante ensayos sobre cilindros curados según \textbf{ASTM C192} (ASTM C192, 2021).

En DIAPCE, la relación a/c opera como un nodo clave del flujo de decisión: el sistema restringe las opciones disponibles a los valores presentes en el conjunto experimental, mostrando únicamente combinaciones observadas en laboratorio (Jiménez y Alvarado, 2020).

% ─────────────────────────────────────────────────────────────────────────────
\subsection{Condiciones ambientales y su efecto sobre el concreto}

El comportamiento del concreto durante su fraguado está determinado en gran medida por las condiciones ambientales de la reacción de hidratación. Las dos variables de mayor incidencia son la \textbf{temperatura} y la \textbf{humedad relativa}, razón por la cual DIAPCE las incorpora como los dos primeros filtros del flujo: temperatura $\rightarrow$ humedad $\rightarrow$ relación a/c $\rightarrow$ aditivo.

\subsubsection*{Temperatura}

La hidratación del cemento es una reacción exotérmica cuya velocidad aumenta con la temperatura. En climas cálidos ---como los que predominan en el Meta, donde las temperaturas medias oscilan entre 25 y 35~°C--- la hidratación se acelera, lo que puede resultar en ganancia rápida de resistencia temprana pero resistencia final menor si no se controlan las condiciones de curado. Temperaturas superiores a 30~°C incrementan la evaporación del agua de mezcla, reducen el tiempo de trabajabilidad y pueden inducir fisuración plástica en las primeras horas.

\subsubsection*{Humedad relativa}

Una humedad relativa baja, combinada con temperatura elevada, puede provocar pérdida acelerada de agua antes de que el concreto alcance suficiente rigidez, generando fisuración por retracción plástica. El curado adecuado ---mediante láminas de polietileno, compuestos de curado o aspersión de agua--- busca contrarrestar esta pérdida (ASTM C192, 2021). En el contexto del Meta, la combinación de temperatura elevada con humedad variable crea condiciones que se desvían de los supuestos implícitos en tablas genéricas de diseño.

% ─────────────────────────────────────────────────────────────────────────────
\subsection{Aditivos químicos para concreto}

Los aditivos químicos son materiales distintos del agua, los agregados y el cemento que se incorporan a la mezcla en pequeñas proporciones con el propósito de modificar alguna de sus propiedades (ASTM C494, 2022). Su uso está regulado por \textbf{ASTM C494/C494M} y por el equivalente colombiano \textbf{NTC 1299}. La selección guiada en DIAPCE muestra únicamente los códigos válidos para la combinación temperatura–humedad–a/c seleccionada.

\subsubsection*{Impermeabilizantes}

Actúan reduciendo la permeabilidad del concreto endurecido mediante la obturación de poros capilares (ASTM C494, 2022). En DIAPCE se representan con los códigos \textbf{P1}, \textbf{P2} y \textbf{P3}, correspondientes a dosificaciones del 2, 3 y 4~\% sobre la masa de cemento del producto \textit{Euco Vandex AM~10I}. El código \textbf{P0} representa la mezcla de control, sin aditivo.

\subsubsection*{Plastificantes}

Permiten mejorar la trabajabilidad sin incrementar la relación a/c (ASTM C494, 2022). En DIAPCE se representan con los códigos \textbf{PP1}, \textbf{PP2} y \textbf{PP3}, correspondientes a dosificaciones del 0{,}2, 0{,}4 y 0{,}6~\% sobre la masa de cemento del producto \textit{Sika\textsuperscript{\textregistered} Plastiment\textsuperscript{\textregistered} AP} (García y Torres, 2021).

\begin{table}[H]
  \centering
  \caption{Clasificación de aditivos en DIAPCE (fuente: datos experimentales del semillero UCC)}
  \begin{tabular}{clll}
    \hline
    \textbf{Código} & \textbf{Tipo} & \textbf{Producto} & \textbf{Dosificación} \\
    \hline
    \textbf{P0}  & Control (sin aditivo) & N/A & 0\,\% \\
    \textbf{P1}  & Impermeabilizante & Euco Vandex AM 10I & 2\,\% \\
    \textbf{P2}  & Impermeabilizante & Euco Vandex AM 10I & 3\,\% \\
    \textbf{P3}  & Impermeabilizante & Euco Vandex AM 10I & 4\,\% \\
    \textbf{PP1} & Plastificante & Sika\textsuperscript{\textregistered} Plastiment\textsuperscript{\textregistered} AP & 0{,}2\,\% \\
    \textbf{PP2} & Plastificante & Sika\textsuperscript{\textregistered} Plastiment\textsuperscript{\textregistered} AP & 0{,}4\,\% \\
    \textbf{PP3} & Plastificante & Sika\textsuperscript{\textregistered} Plastiment\textsuperscript{\textregistered} AP & 0{,}6\,\% \\
    \hline
  \end{tabular}
\end{table}

% ─────────────────────────────────────────────────────────────────────────────
\subsection{Desarrollo de resistencia a compresión}

La resistencia a compresión no se alcanza de forma instantánea: es el resultado de un proceso continuo de hidratación que se extiende durante semanas o meses (ASTM C39, 2021). Por convención, la resistencia de diseño $f'_c$ se especifica a los \textbf{28 días}. Las edades intermedias de \textbf{7 días} y \textbf{14 días} tienen valor diagnóstico (ASTM C192, 2021). DIAPCE estima las resistencias como promedios empíricos del subconjunto filtrado:
\[
\hat{f}_{t} \;=\; \operatorname{avg}\bigl\{\, f_{t,i} \;:\; (T_i,\, H_i,\, \mathrm{a/c}_i,\, \mathrm{aditivo}_i) \in \text{selección}\,\bigr\},
\quad t \in \{7,\, 14,\, 28\}.
\]

% ─────────────────────────────────────────────────────────────────────────────
\subsection{Modelado empírico y enfoques de aprendizaje}

DIAPCE adopta un \textbf{modelo empírico de agregación} (promedios condicionados) por su transparencia, robustez ante datos ruidosos y compatibilidad con las decisiones en cascada (Jiménez y Alvarado, 2020; Torres y Jiménez, 2023). La arquitectura deja abierta la integración futura de modelos ML ---como regresión regularizada o \textit{tree ensembles}--- entrenados sobre las mismas vistas SQL.

% ─────────────────────────────────────────────────────────────────────────────
\subsection{Fundamentos tecnológicos de la solución}

\subsubsection*{Flutter y Dart}
Flutter es un \textit{framework} de código abierto de Google que permite construir aplicaciones nativas para Android, iOS, web y escritorio desde una única base de código en Dart (Google LLC, 2024). La interfaz de DIAPCE implementa formularios con validadores y visualizaciones con \texttt{fl\_chart}.

\subsubsection*{SQLite y almacenamiento local}
SQLite opera directamente sobre el dispositivo sin requerir servidor externo (SQLite Consortium, 2024). En DIAPCE gestiona parámetros de diseño, resultados de laboratorio y datos experimentales del CSV, con normalización, llaves foráneas e índices para consultas condicionales, vía la biblioteca \texttt{sqflite} (Tekartik, 2024).

\subsubsection*{Arquitectura limpia y datos en CSV}
El dataset se almacena en \texttt{csv\_datos\_1.1.csv} y se inserta en SQLite al iniciar la aplicación por primera vez. La arquitectura sigue \textit{Clean Architecture}: dominio, datos y presentación.

% ─────────────────────────────────────────────────────────────────────────────
\subsection{Trabajos relacionados}

\textbf{ConcreBrain} utiliza redes neuronales para predecir resistencia; \textbf{Concrete Mix Design} es una solución comercial estándar; \textbf{Concretelab} facilita diseño en entorno web. DIAPCE se diferencia por operar completamente sin conexión, derivar sus recomendaciones del dataset propio del semillero de la UCC y tener un diseño explícitamente formativo e institucional. Como línea base normativa: ACI 211.1 (2016), ASTM C39 (2021), ASTM C192 (2021), ASTM C494 (2022).

% ══════════════════════════════════════════════════════════════════════════════
%  CAPÍTULO 3 — DISEÑO METODOLÓGICO
% ══════════════════════════════════════════════════════════════════════════════
\newpage
\section{Diseño metodológico}

Este capítulo describe el proceso de diseño y construcción de DIAPCE, organizado en cinco fases: levantamiento de requisitos, modelado de datos, diseño de la solución, desarrollo e implementación, y pruebas de calidad.

% ─────────────────────────────────────────────────────────────────────────────
\subsection{Requisitos}

Los requisitos de DIAPCE se identificaron a partir de las necesidades del semillero de investigación y de los laboratorios de materiales de la UCC.

\subsubsection*{Requisitos funcionales}
\begin{enumerate}
    \item Autenticación local de usuarios (registro e inicio de sesión).
    \item Crear, visualizar y eliminar proyectos de diseño de mezclas.
    \item Flujo de selección en cascada: resistencia objetivo $\rightarrow$ temperatura $\rightarrow$ humedad $\rightarrow$ relación a/c $\rightarrow$ aditivo (P0--PP3).
    \item Calcular y mostrar proyecciones de resistencia a 7, 14 y 28 días basadas en el dataset experimental.
    \item Almacenar de forma persistente todos los parámetros de entrada y las predicciones generadas.
    \item Operar sin conexión a internet (\textit{offline-first}).
    \item Mostrar gráficos de composición de mezcla (pastel) y evolución de resistencia (línea).
\end{enumerate}

\subsubsection*{Requisitos no funcionales}
\begin{itemize}
    \item \textbf{Portabilidad:} Android, iOS, web y escritorio desde una única base de código.
    \item \textbf{Rendimiento:} consultas de predicción en menos de 500\,ms en dispositivos de gama media.
    \item \textbf{Trazabilidad:} cada proyecto registra parámetros de entrada, versión del dataset y fecha de creación.
    \item \textbf{Usabilidad:} interfaz intuitiva para estudiantes sin experiencia previa.
    \item \textbf{Mantenibilidad:} arquitectura que permite incorporar nuevos modelos sin reescribir la capa de presentación.
\end{itemize}

% ─────────────────────────────────────────────────────────────────────────────
\subsection{Modelado}

\subsubsection*{Esquema de la base de datos del servidor (PostgreSQL)}

Como parte del análisis de los datos experimentales, se diseñó un esquema normalizado en PostgreSQL. El proceso de normalización progresiva resolvió problemas de ambigüedad y redundancia, derivando en un diseño final de cuatro tablas interrelacionadas.

\begin{lstlisting}[style=dbmlstyle, caption=Esquema de la base de datos en DBML (PostgreSQL).]
Table TiposAditivo {
  id integer [primary key, increment]
  nombre varchar(50) [unique, not null]
}
Table Productos {
  id integer [primary key, increment]
  nombre_producto varchar(100) [unique, not null]
  marca varchar(50)
}
Table Aditivos {
  id integer [primary key, increment]
  codigo varchar(10) [unique, not null]
  porcentaje_aplicado varchar(10) [not null]
  tipo_aditivo_id integer [not null, ref: > TiposAditivo.id]
  producto_id integer [not null, ref: > Productos.id]
}
Table ResultadosConcreto {
  id integer [primary key, increment]
  temperatura integer [not null]
  humedad integer [not null]
  relacion_ac numeric(3,2) [not null]
  edad_dias integer [not null]
  resistencia_mpa numeric(10,2) [not null]
  aditivo_id integer [not null, ref: > Aditivos.id]
}
\end{lstlisting}

\begin{figure}[H]
    \centering
    \includegraphics[width=1\textwidth]{images/diagrama_db.png}
    \caption{Diagrama Entidad-Relación de la estructura final de 4 tablas.}
    \label{fig:er-postgres}
\end{figure}

\subsubsection*{Esquema de la base de datos local (SQLite)}

La base de datos local extiende el esquema anterior con tablas propias de la aplicación: \texttt{users}, \texttt{projects} y \texttt{mixtures}. Un índice compuesto en \texttt{ResultadosConcreto} optimiza las consultas de predicción en cascada.

\begin{figure}[H]
    \centering
    \includegraphics[width=1\textwidth]{images/diagrama_db_completo.jpg}
    \caption{Diagrama Entidad-Relación de la base de datos local completa (SQLite).}
    \label{fig:er-sqlite}
\end{figure}

El siguiente código DBML detalla la estructura completa de la base de datos local:

\begin{lstlisting}[style=dbmlstyle, caption=Esquema completo de la base de datos local en DBML.]
Table users {
  id integer [pk, increment]
  email string [unique, not null]
  password string [not null]
}
Table projects {
  id integer [pk, increment]
  user_id integer [not null, ref: > users.id]
  project_name string [not null]
  selected_date string
  selected_image_path string
  creator_name string
  resistance_target integer [not null]
  temperature integer [not null]
  humidity integer [not null]
  relacion_ac float [not null]
  work_type string
  aditivo_id integer [ref: > Aditivos.id]
  resistencia_predicha_7d float
  resistencia_predicha_14d float
  resistencia_predicha_28d float
  created_at string
}
Table TiposAditivo {
  id integer [pk, increment]
  nombre string [unique, not null]
}
Table Productos {
  id integer [pk, increment]
  nombre_producto string [unique, not null]
  marca string
}
Table Aditivos {
  id integer [pk, increment]
  codigo string [unique, not null]
  porcentaje_aplicado string [not null]
  tipo_aditivo_id integer [not null, ref: > TiposAditivo.id]
  producto_id integer [not null, ref: > Productos.id]
}
Table ResultadosConcreto {
  id integer [pk, increment]
  temperatura integer [not null]
  humedad integer [not null]
  relacion_ac float [not null]
  edad_dias integer [not null]
  resistencia_mpa float [not null]
  aditivo_id integer [not null, ref: > Aditivos.id]
  indexes {
    (temperatura, humedad, relacion_ac, aditivo_id, edad_dias)
      [name: 'idx_busqueda_resistencia']
  }
}
Table mixtures {
  id integer [pk, increment]
  project_id integer [ref: > projects.id]
  name string
  description string
  cemento_kg float
  agua_litros float
  arena_kg float
  grava_kg float
  aditivo_porcentaje float
  created_at string
}
\end{lstlisting}

% ─────────────────────────────────────────────────────────────────────────────
\subsection{Diseño}

\subsubsection*{Planteamiento de la solución: DIAPCE}

DIAPCE es una aplicación desarrollada en Flutter, diseñada como herramienta de apoyo para el diseño de mezclas de concreto, fundamentada en datos experimentales locales del semillero de la UCC. Su objetivo es asistir a los estudiantes en la selección de mezclas basando las recomendaciones en datos reales en lugar de estimaciones externas.

\subsubsection*{Propósito y flujo principal}

La aplicación permite crear y gestionar proyectos donde se definen condiciones controladas y se obtienen predicciones de resistencia a los 7, 14 y 28 días. El flujo de usuario es el siguiente:

\begin{enumerate}
    \item El usuario inicia la aplicación e ingresa a su perfil.
    \item Crea un nuevo proyecto especificando nombre, fecha, tipo de obra e imagen opcional.
    \item Define las propiedades técnicas mediante selección en cascada: resistencia objetivo, temperatura, humedad, relación a/c y aditivo.
    \item La aplicación consulta el dataset, calcula los promedios de resistencia y presenta una predicción gráfica y numérica.
    \item El proyecto, junto con sus condiciones y predicciones, se guarda en SQLite asociado al perfil del usuario.
\end{enumerate}

\subsubsection*{Arquitectura del sistema}

DIAPCE sigue el enfoque de \textit{Clean Architecture}, separando responsabilidades en tres capas:
\begin{itemize}
    \item \textbf{Dominio:} entidades de negocio (Proyecto, Mezcla, ResultadoConcreto) y casos de uso.
    \item \textbf{Datos:} fuentes locales (SQLite vía \texttt{sqflite}), repositorios y modelos de mapeo.
    \item \textbf{Presentación:} pantallas Flutter, \textit{widgets} reutilizables y gestión de estado.
\end{itemize}

\subsubsection*{Flujo de decisión en cascada}

\[
\text{Resistencia objetivo} \;\rightarrow\; \text{Temperatura} \;\rightarrow\; \text{Humedad} \;\rightarrow\; \text{Relación a/c} \;\rightarrow\; \text{Aditivo (P0--PP3)}
\]
En cada paso, el sistema muestra únicamente los valores disponibles para las condiciones ya seleccionadas, eliminando combinaciones sin evidencia experimental.

% ─────────────────────────────────────────────────────────────────────────────
\subsection{Desarrollo}

\subsubsection*{Tecnologías utilizadas}
\begin{itemize}
    \item \textbf{Flutter / Dart:} framework multiplataforma para interfaz y lógica de la aplicación.
    \item \textbf{sqflite:} biblioteca para la gestión de SQLite local.
    \item \textbf{fl\_chart:} biblioteca para gráficos interactivos (pastel y línea).
    \item \textbf{csv:} biblioteca para importar el dataset experimental desde \texttt{csv\_datos\_1.1.csv}.
\end{itemize}

\subsubsection*{Implementación en PostgreSQL}

\begin{lstlisting}[style=sqlstyle, caption=Script DDL para PostgreSQL.]
CREATE TABLE TiposAditivo (
    id SERIAL PRIMARY KEY,
    nombre VARCHAR(50) UNIQUE NOT NULL
);
CREATE TABLE Productos (
    id SERIAL PRIMARY KEY,
    nombre_producto VARCHAR(100) UNIQUE NOT NULL,
    marca VARCHAR(50)
);
CREATE TABLE Aditivos (
    id SERIAL PRIMARY KEY,
    codigo VARCHAR(10) UNIQUE NOT NULL,
    porcentaje_aplicado VARCHAR(10) NOT NULL,
    tipo_aditivo_id INT NOT NULL REFERENCES TiposAditivo(id),
    producto_id INT NOT NULL REFERENCES Productos(id)
);
CREATE TABLE ResultadosConcreto (
    id SERIAL PRIMARY KEY,
    temperatura INT NOT NULL,
    humedad INT NOT NULL,
    relacion_ac DECIMAL(3, 2) NOT NULL,
    edad_dias INT NOT NULL,
    resistencia_mpa DECIMAL(10, 2) NOT NULL,
    aditivo_id INT NOT NULL REFERENCES Aditivos(id)
);
\end{lstlisting}

\begin{lstlisting}[style=sqlstyle, caption=Cálculo de la resistencia promedio por grupo experimental.]
SELECT
    p.nombre_producto,
    a.codigo AS codigo_aditivo,
    r.edad_dias,
    ROUND(AVG(r.resistencia_mpa), 2) AS promedio_resistencia,
    COUNT(r.id) AS numero_de_muestras
FROM ResultadosConcreto AS r
JOIN Aditivos AS a ON r.aditivo_id = a.id
JOIN Productos AS p ON a.producto_id = p.id
GROUP BY p.nombre_producto, a.codigo, r.edad_dias
ORDER BY a.codigo, r.edad_dias;
\end{lstlisting}

\subsubsection*{Flujo de datos y consulta clave de predicción}

Al inicializar, la aplicación carga los datos desde el CSV a \texttt{ResultadosConcreto}. Esta copia local permite calcular predicciones de manera instantánea, sin depender de una conexión constante.

\begin{figure}[H]
    \centering
    \includegraphics[width=0.8\textwidth]{images/flujo_prediccion.jpg}
    \caption{Flujo de datos para el cálculo de predicciones en la base de datos local.}
    \label{fig:prediction_flow}
\end{figure}

\begin{lstlisting}[style=sqlstyle, caption=Consulta clave para calcular promedios de resistencia.]
SELECT
    edad_dias,
    ROUND(AVG(resistencia_mpa), 2) AS resistencia_promedio,
    COUNT(*) AS num_muestras
FROM resultados_concreto
WHERE temperatura = ? AND humedad = ? AND relacion_ac = ? AND aditivo_id = ?
GROUP BY edad_dias
ORDER BY edad_dias ASC;
\end{lstlisting}

\subsubsection*{Vistas de la aplicación}

\paragraph{Inicio de sesión (\texttt{lib/view/main\_login.dart}).}
Autenticación local de usuarios con validación básica de campos y alternancia de visibilidad de contraseña.

\begin{figure}[H]
  \centering
  \includegraphics[width=0.45\linewidth]{images/main_login.jpg}
  \caption{Pantalla de inicio de sesión de la aplicación.}
  \label{fig:view-login}
\end{figure}

\paragraph{Crear cuenta (\texttt{lib/view/create\_count.dart}).}
Registro de nuevos usuarios con confirmación de contraseña y validaciones de formato.

\begin{figure}[H]
  \centering
  \includegraphics[width=0.45\linewidth]{images/create_count.jpg}
  \caption{Pantalla de registro de nuevos usuarios.}
  \label{fig:view-registro}
\end{figure}

\paragraph{Panel de proyectos --- Hall (\texttt{lib/view/hall.dart}).}
Listado en cuadrícula de los proyectos del usuario con acciones de creación y eliminación con confirmación.

\begin{figure}[H]
  \centering
  \includegraphics[width=0.45\linewidth]{images/hall_vacio.jpg}
  \caption{Panel principal (Hall) sin proyectos.}
  \label{fig:view-hall-vacio}
\end{figure}

\begin{figure}[H]
  \centering
  \includegraphics[width=0.45\linewidth]{images/hall_lleno.jpg}
  \caption{Panel principal (Hall) con proyectos del usuario.}
  \label{fig:view-hall-lleno}
\end{figure}

\paragraph{Crear proyecto (\texttt{lib/view/create\_proyect\_screen.dart}).}
Flujo de selección en cascada con campo de resistencia objetivo validado en el rango [24{,}00,\,57{,}00]\,MPa.

\begin{figure}[H]
  \centering
  \includegraphics[width=0.45\linewidth]{images/create_proyect_screen.jpg}
  \caption{Pantalla de creación de proyecto con flujo en cascada (parte 1).}
  \label{fig:view-crear-proyecto-1}
\end{figure}

\begin{figure}[H]
  \centering
  \includegraphics[width=0.45\linewidth]{images/create_proyect_screen_2.jpg}
  \caption{Pantalla de creación de proyecto con flujo en cascada (parte 2).}
  \label{fig:view-crear-proyecto-2}
\end{figure}

\paragraph{Detalle de proyecto (\texttt{lib/view/ViewExistingProjectScreen.dart}).}
Visualización integral con predicciones a 7/14/28 días, gráfico de composición (pastel) y evolución de resistencia (línea).

\begin{figure}[H]
  \centering
  \includegraphics[width=0.45\linewidth]{images/ViewExistingProjectScreen.jpg}
  \caption{Detalle de un proyecto existente con gráficos de predicción (parte 1).}
  \label{fig:view-detalle-1}
\end{figure}

\begin{figure}[H]
  \centering
  \includegraphics[width=0.45\linewidth]{images/ViewExistingProjectScreen_2.jpg}
  \caption{Detalle de un proyecto existente con gráficos de predicción (parte 2).}
  \label{fig:view-detalle-2}
\end{figure}

\begin{figure}[H]
  \centering
  \includegraphics[width=0.45\linewidth]{images/ViewExistingProjectScreen_3.jpg}
  \caption{Detalle de un proyecto existente con gráficos de predicción (parte 3).}
  \label{fig:view-detalle-3}
\end{figure}

\paragraph{Vista de propiedades --- legado (\texttt{lib/view/view\_properties.dart}).}
Pantalla histórica de propiedades; actualmente desactivada/comentada. Útil como referencia de diseño.

\begin{figure}[H]
  \centering
  \includegraphics[width=0.45\linewidth]{images/view_properties.jpg}
  \caption{Vista de propiedades (legado).}
  \label{fig:view-propiedades}
\end{figure}

% ─────────────────────────────────────────────────────────────────────────────
\subsection{Pruebas}

Las pruebas se organizan en tres niveles: funcionales, de calidad del modelo y de cumplimiento normativo.

\subsubsection*{Pruebas funcionales}
\begin{itemize}
    \item Registro e inicio de sesión con credenciales válidas e inválidas.
    \item Creación de proyecto con todos los campos del flujo en cascada completos.
    \item Eliminación de proyecto con confirmación y verificación de borrado en cascada.
    \item Visualización correcta de gráficos tras recuperar una predicción del dataset.
    \item Operación completa sin conexión a internet en dispositivo Android.
\end{itemize}

\subsubsection*{Pruebas de calidad del modelo de predicción}

La validación se realizó mediante \textit{Leave-One-Out Cross-Validation} (LOO) sobre los 1\,872 registros del dataset. Los resultados se detallan en el Capítulo~4.

\subsubsection*{Pruebas de cumplimiento normativo}

Se verificó que las recomendaciones sean coherentes con los rangos establecidos por ACI\,211.1 para los niveles de temperatura (10, 25 y 32~°C), humedad (20, 50 y 70~\%) y relación a/c (0{,}40, 0{,}45 y 0{,}50).

% ══════════════════════════════════════════════════════════════════════════════
%  CAPÍTULO 4 — RESULTADOS, LIMITACIONES Y LOGROS
% ══════════════════════════════════════════════════════════════════════════════
\newpage
\section{Resultados, limitaciones y logros}

Esta sección analiza los datos y resultados generados con DIAPCE en el contexto del Meta (Colombia), con énfasis en su uso formativo por parte de \textbf{los estudiantes del semillero de investigación y de los laboratorios de materiales de la UCC}. La discusión se centra en: (i) la calidad y cobertura del conjunto de datos regional; (ii) la validez de las recomendaciones del sistema para dosificación base y clasificación de aditivos (P0--PP3) con proyecciones a 7, 14 y 28 días; y (iii) la comparación frente a prácticas tradicionales de diseño según ACI\,211.1.

% ─────────────────────────────────────────────────────────────────────────────
\subsection{Conjunto de datos y preprocesamiento}

El conjunto de datos registra parámetros ambientales y de diseño (temperatura, humedad, relación a/c, condición de exposición) y variables de mezcla y desempeño.

\begin{table}[H]
  \centering
  \caption{Cobertura de rangos del conjunto de datos (1\,872 registros)}
  \begin{tabular}{lccp{5.5cm}}
    \hline
    Variable & Mín & Máx & Observaciones \\
    \hline
    Resistencia observada (MPa) & 22{,}57 & 58{,}96 & P$_{25}$=31{,}19 · P$_{50}$=37{,}97 · P$_{75}$=43{,}34 MPa \\
    Temperatura (°C) & 10 & 32 & 3 niveles: 10\,°C (30{,}8\%), 25\,°C (46{,}2\%), 32\,°C (23{,}1\%) \\
    Humedad relativa (\%) & 20 & 70 & 3 niveles: 20\% (23{,}1\%), 50\% (46{,}2\%), 70\% (30{,}8\%) \\
    Relación a/c & 0{,}40 & 0{,}50 & 3 niveles: 0{,}40 (30{,}8\%), 0{,}45 (30{,}8\%), 0{,}50 (38{,}5\%) \\
    \hline
  \end{tabular}
\end{table}

% ─────────────────────────────────────────────────────────────────────────────
\subsection{Metodología de validación}

Se empleó \textit{Leave-One-Out Cross-Validation} (LOO) sobre el dataset completo. Las métricas consideradas son:
\begin{itemize}
  \item \textbf{Clasificación de aditivos} (P0--PP3): exactitud, precisión, exhaustividad y F1 macro.
  \item \textbf{Proyección de resistencia} (7/14/28 días): MAE, RMSE y MAPE por horizonte.
  \item \textbf{Coherencia normativa:} cumplimiento de rangos indicativos según ACI\,211.1.
  \item \textbf{Trazabilidad:} porcentaje de recomendaciones con insumos y versionado correctamente almacenados en SQLite.
\end{itemize}

% ─────────────────────────────────────────────────────────────────────────────
\subsection{Resultados de clasificación de aditivos}

\begin{table}[H]
  \centering
  \caption{Métricas de clasificación de aditivos --- vecino más cercano LOO (1\,872 registros, 7 clases)}
  \begin{tabular}{lcccc}
    \hline
    Clase & Precisión & Exhaustividad & F1 & Soporte \\
    \hline
    P0  & 0{,}537 & 0{,}528 & 0{,}532 & 468 \\
    P1  & 0{,}187 & 0{,}192 & 0{,}189 & 234 \\
    P2  & 0{,}156 & 0{,}132 & 0{,}143 & 234 \\
    P3  & 0{,}154 & 0{,}137 & 0{,}145 & 234 \\
    PP1 & 0{,}128 & 0{,}124 & 0{,}126 & 234 \\
    PP2 & 0{,}164 & 0{,}158 & 0{,}161 & 234 \\
    PP3 & 0{,}383 & 0{,}513 & 0{,}439 & 234 \\
    \textbf{Macro-promedio} & 0{,}244 & 0{,}255 & 0{,}248 & 1\,872 \\
    \hline
  \end{tabular}
\end{table}

La clase P0 (control, sin aditivo) mostró el mejor desempeño individual (F1\,=\,0{,}532). Esto confirma que la \textbf{selección guiada en cascada} es el enfoque correcto, ya que la resistencia sola no discrimina suficientemente entre aditivos.

% ─────────────────────────────────────────────────────────────────────────────
\subsection{Resultados de proyección de resistencia}

\begin{table}[H]
  \centering
  \caption{Error de proyección de resistencia --- validación LOO por horizonte}
  \begin{tabular}{lccc}
    \hline
    Horizonte & MAE (MPa) & RMSE (MPa) & MAPE (\%) \\
    \hline
    7 días  & 1{,}284 & 1{,}625 & 4{,}40 \\
    14 días & 1{,}767 & 2{,}207 & 4{,}61 \\
    28 días & 1{,}965 & 2{,}455 & 4{,}37 \\
    \hline
  \end{tabular}
\end{table}

Los errores de proyección son menores a 2\,MPa en todos los horizontes evaluados, con MAPE inferior al 5\,\%, adecuado para uso formativo.

% ─────────────────────────────────────────────────────────────────────────────
\subsection{Comparación con métodos tradicionales}

Para efectos académicos, se compararon las recomendaciones de DIAPCE con diseños elaborados mediante procedimiento manual guiado por ACI\,211.1. Las recomendaciones generadas mostraron coherencia con los rangos indicativos de ACI\,211.1 en las combinaciones cubiertas por el dataset.

\begin{figure}[H]
  \centering
  \fbox{\rule{0pt}{2in}\rule{0.9\linewidth}{0pt}}
  \caption{Comparación de resultados: DIAPCE vs. diseño manual ACI (espacio reservado para figura).}
  \label{fig:comparacion-aci}
\end{figure}

% ─────────────────────────────────────────────────────────────────────────────
\subsection{Sostenibilidad e impacto}

Aunque DIAPCE no implementa optimización automática de proporciones, su uso sistemático puede incentivar decisiones de dosificación más consistentes. Cuando se disponga de mediciones, puede evaluarse el impacto potencial en: (i) contenido de cemento frente a una línea base docente, y (ii) variación de resistencia a 28 días.

% ─────────────────────────────────────────────────────────────────────────────
\subsection{Amenazas a la validez y limitaciones}

\begin{itemize}
  \item \textbf{Sesgo de datos:} el conjunto es regional; la generalización fuera del Meta debe justificarse.
  \item \textbf{Medición:} variabilidad experimental en laboratorio puede introducir ruido en las proyecciones.
  \item \textbf{Cobertura:} desempeño reducido en combinaciones poco representadas (extremos de temperatura/humedad).
  \item \textbf{Alcance del sistema:} la herramienta asiste decisiones; no sustituye criterio profesional ni verificación normativa.
  \item El dataset cubre únicamente 3 niveles discretos de temperatura, 3 de humedad y 3 valores de a/c, sin interpolación.
  \item Desbalance de clase: P0 concentra el 25~\% de los registros frente al 12{,}5~\% de cada aditivo restante.
  \item No se dispone de módulo de optimización automática ni de sincronización cliente-servidor.
\end{itemize}

% ─────────────────────────────────────────────────────────────────────────────
\subsection{Reproducibilidad}

DIAPCE almacena localmente entradas y resultados (SQLite), lo que permite rastrear cada recomendación respecto a sus insumos y la versión del modelo. Se sugiere anexar tablas exportables (CSV) con métricas por cohorte de ensayos.

% ─────────────────────────────────────────────────────────────────────────────
\subsection{Logros}

\begin{itemize}
  \item Aplicación multiplataforma funcional, operativa sin conexión, con flujo en cascada validado sobre 1\,872 registros experimentales.
  \item Errores de proyección por debajo de 2\,MPa (MAPE $<$ 5\,\%) en los tres horizontes, adecuados para uso formativo.
  \item Base de datos normalizada con integridad referencial, índices optimizados y trazabilidad completa.
  \item Interfaz intuitiva validada en escenarios de laboratorio con estudiantes del semillero de la UCC.
  \item Diseño offline-first que reduce fricciones de despliegue en entornos con conectividad limitada.
\end{itemize}

% ══════════════════════════════════════════════════════════════════════════════
%  CAPÍTULO 5 — CONCLUSIONES Y RECOMENDACIONES
% ══════════════════════════════════════════════════════════════════════════════
\newpage
\section{Conclusiones y recomendaciones}

% ─────────────────────────────────────────────────────────────────────────────
\subsection{Conclusiones}

Con base en los objetivos planteados y los resultados obtenidos sobre el conjunto de datos regional (1\,872 registros, 24--57\,MPa, tres niveles de temperatura: 10, 25 y 32\,°C; tres niveles de humedad: 20, 50 y 70\,\%; relaciones a/c de 0{,}40, 0{,}45 y 0{,}50), se concluye lo siguiente:

\begin{itemize}
  \item \textbf{Implementación de la herramienta:} Se desarrolló DIAPCE como una aplicación multiplataforma en Flutter con almacenamiento local mediante SQLite, \textbf{orientada a los estudiantes del semillero de investigación y de laboratorio de materiales de la UCC, Meta}. La herramienta asiste la toma de decisiones en el diseño de mezclas: registra parámetros de entrada, sugiere una dosificación base y clasifica el uso de aditivos (\textbf{P0}: sin aditivo; \textbf{P1}--\textbf{P3}: impermeabilizante al 2, 3 y~4\,\%; \textbf{PP1}--\textbf{PP3}: plastificante al 0{,}2, 0{,}4 y~0{,}6\,\%), entregando proyecciones de resistencia a 7, 14 y 28 días en concordancia con ACI~211.1. Su alcance es institucional y académico; no está diseñada para uso industrial ni general.

  \item \textbf{Enfoque estudiantil y trazabilidad:} Se consolidó un flujo de trabajo pensado desde y para los estudiantes de los semilleros y laboratorios de la UCC, enfatizando trazabilidad (registro de insumos, resultados y versión del modelo) para favorecer reproducibilidad, evaluación docente y comparación entre ensayos.

  \item \textbf{Desempeño en proyección de resistencia:} La validación Leave-One-Out arrojó errores de proyección menores a 2\,MPa en todos los horizontes: MAE\,=\,1{,}28\,MPa, RMSE\,=\,1{,}63\,MPa y MAPE\,=\,4{,}40\,\% a 7 días; MAE\,=\,1{,}77\,MPa, RMSE\,=\,2{,}21\,MPa y MAPE\,=\,4{,}61\,\% a 14 días; MAE\,=\,1{,}97\,MPa, RMSE\,=\,2{,}46\,MPa y MAPE\,=\,4{,}37\,\% a 28 días.

  \item \textbf{Clasificación guiada de aditivos:} La clasificación automática por vecino más cercano (LOO) obtuvo un macro-F1\,=\,0{,}248. Esto confirma que la selección guiada en cascada ---y no la clasificación automática por resistencia--- es el enfoque correcto para la herramienta.

  \item \textbf{Adopción y operación:} El diseño offline-first y multiplataforma reduce fricciones de despliegue en entornos con conectividad limitada, facilitando su adopción en aula y laboratorio.

  \item \textbf{Contribución y límites:} DIAPCE no sustituye el criterio del docente ni el cumplimiento normativo; lo complementa con evidencia sistematizada. Las limitaciones principales son la cobertura discreta del dataset y la ausencia de sincronización cliente-servidor, que definen líneas claras de trabajo futuro.
\end{itemize}

% ─────────────────────────────────────────────────────────────────────────────
\subsection{Recomendaciones}

\begin{itemize}
  \item \textbf{Ampliar y balancear el conjunto de datos:} Incorporar más ensayos para los aditivos P1--P3 y PP1--PP3, ampliar los niveles de temperatura y humedad, y extender el rango de a/c más allá de los tres valores actuales.

  \item \textbf{Mejorar la discriminación entre aditivos:} Explorar variables adicionales de discriminación (trabajabilidad, tiempo de fraguado, contenido de cemento) que permitan separar mejor las clases.

  \item \textbf{Explorar modelos supervisados (ML):} Modelos como regresión regularizada o \textit{tree ensembles} entrenados sobre las mismas vistas SQL podrían reducir el error de proyección en rangos menos representados del dataset.

  \item \textbf{Evaluación continua y exportación de métricas:} Establecer un flujo reproducible que calcule y exporte periódicamente (CSV/\LaTeX) las métricas: F1 por clase (P0--PP3), MAE/RMSE/MAPE por horizonte (7/14/28 días) y coherencia ACI.

  \item \textbf{Validación con ensayos reales de laboratorio:} Contrastar las proyecciones de DIAPCE con las resistencias medidas experimentalmente por los estudiantes (ASTM C39, NTC 673).

  \item \textbf{Usabilidad y docencia:} Incorporar exportación a PDF/CSV de proyectos y resultados, plantillas de prácticas adaptadas a los cursos de la UCC y tutoriales integrados.

  \item \textbf{Escalabilidad opcional:} Evaluar a futuro una capa de sincronización/backend para consolidar datos entre dispositivos, preservando siempre la operación offline como principio de diseño.

  \item \textbf{Gobernanza de datos:} Mantener versionado de datasets y modelos, control de calidad y lineamientos éticos para el uso de información experimental generada por los estudiantes.

  \item \textbf{Generalización regional (trabajo futuro):} Una vez consolidado el dataset del Meta, evaluar la adaptación del sistema para otras regiones de Colombia, manteniendo la validación experimental local como requisito previo.
\end{itemize}

% ══════════════════════════════════════════════════════════════════════════════
%  REFERENCIAS BIBLIOGRÁFICAS
% ══════════════════════════════════════════════════════════════════════════════
\newpage
\addcontentsline{toc}{section}{Referencias bibliográficas}
\printbibliography[title={Referencias bibliográficas}]

% ══════════════════════════════════════════════════════════════════════════════
%  ANEXOS
% ══════════════════════════════════════════════════════════════════════════════
\newpage
\section*{Anexos}
\addcontentsline{toc}{section}{Anexos}

\subsection*{Anexo A. Estructura del dataset experimental (\texttt{csv\_datos\_1.1.csv})}

\begin{table}[H]
  \centering
  \caption{Campos del dataset experimental}
  \begin{tabular}{lll}
    \hline
    Campo & Tipo & Descripción \\
    \hline
    \texttt{temperatura}      & Entero  & Temperatura de curado (°C): 10, 25 o 32 \\
    \texttt{humedad}          & Entero  & Humedad relativa (\%): 20, 50 o 70 \\
    \texttt{relacion\_ac}     & Decimal & Relación agua/cemento: 0{,}40; 0{,}45 o 0{,}50 \\
    \texttt{edad\_dias}       & Entero  & Edad del ensayo: 7, 14 o 28 días \\
    \texttt{resistencia\_mpa} & Decimal & Resistencia medida (MPa) \\
    \texttt{codigo\_aditivo}  & Texto   & Código del aditivo: P0, P1, P2, P3, PP1, PP2 o PP3 \\
    \hline
  \end{tabular}
\end{table}

\subsection*{Anexo B. Script SQL de creación de tablas (PostgreSQL)}

\begin{lstlisting}[style=sqlstyle, caption=Script DDL completo para PostgreSQL.]
CREATE TABLE TiposAditivo (
    id SERIAL PRIMARY KEY,
    nombre VARCHAR(50) UNIQUE NOT NULL
);
CREATE TABLE Productos (
    id SERIAL PRIMARY KEY,
    nombre_producto VARCHAR(100) UNIQUE NOT NULL,
    marca VARCHAR(50)
);
CREATE TABLE Aditivos (
    id SERIAL PRIMARY KEY,
    codigo VARCHAR(10) UNIQUE NOT NULL,
    porcentaje_aplicado VARCHAR(10) NOT NULL,
    tipo_aditivo_id INT NOT NULL REFERENCES TiposAditivo(id),
    producto_id INT NOT NULL REFERENCES Productos(id)
);
CREATE TABLE ResultadosConcreto (
    id SERIAL PRIMARY KEY,
    temperatura INT NOT NULL,
    humedad INT NOT NULL,
    relacion_ac DECIMAL(3, 2) NOT NULL,
    edad_dias INT NOT NULL,
    resistencia_mpa DECIMAL(10, 2) NOT NULL,
    aditivo_id INT NOT NULL REFERENCES Aditivos(id)
);
\end{lstlisting}

\end{document}
